% ============================================================
% Introduzione
% ============================================================
\chapter*{Introduzione}
\addcontentsline{toc}{chapter}{Introduzione}

% --- Contesto generale ---
La trasformazione digitale della Pubblica Amministrazione rappresenta una delle sfide strategiche più rilevanti nel panorama europeo e italiano. L'adozione di tecnologie basate sull'Intelligenza Artificiale (IA) si inserisce in questo processo come leva fondamentale per il miglioramento dell'efficienza, della qualità dei servizi e della capacità decisionale delle organizzazioni pubbliche.

% --- Problema di ricerca ---
In tale contesto, la scelta del modello di deployment --- cloud, on-premise o ibrido --- assume un ruolo cruciale, soprattutto in relazione ai requisiti di sovranità digitale, protezione dei dati, compliance normativa e continuità operativa. Il paradigma on-premise, pur presentando sfide in termini di costi e competenze, offre vantaggi significativi per le amministrazioni che trattano dati sensibili e necessitano di pieno controllo sulle infrastrutture tecnologiche.

% --- Obiettivo della tesi ---
Il presente lavoro si propone di analizzare l'utilizzo di sistemi di Intelligenza Artificiale in ambiente on-premise a supporto dei processi decisionali strategici nella Pubblica Amministrazione. L'analisi si sviluppa attraverso la ricostruzione del quadro normativo e strategico europeo e nazionale, l'esame delle linee guida e dei framework di riferimento, e lo studio di casi istituzionali e accademici che hanno adottato soluzioni di IA on-premise.

% --- Struttura della tesi ---
La tesi è articolata in quattro capitoli.

Il \textbf{Capitolo~1} ricostruisce l'evoluzione del paradigma di Digital Government, analizzando il Piano Triennale ICT 2024--2026, le infrastrutture digitali e la cloud strategy nella PA, la cybersecurity e la sovranità digitale, nonché gli indicatori europei di maturità digitale (DESI, Digital Government Index).

Il \textbf{Capitolo~2} affronta il tema dell'Intelligenza Artificiale nella Pubblica Amministrazione, esaminando la trasformazione organizzativa indotta dall'IA, il quadro regolatorio europeo (AI Act, Data Governance Act, Data Act), le Linee guida AgID, i modelli di deployment e le questioni di compliance e accountability algoritmica.

Il \textbf{Capitolo~3} è dedicato alla Data Governance e alla maturità digitale negli Atenei, con riferimento ai framework internazionali (DAMA, ISO, NIST), ai maturity model per la PA e al caso specifico dell'Università degli Studi di Napoli ``Parthenope'', di cui si propone una gap analysis.

Il \textbf{Capitolo~4} presenta una rassegna di casi europei e istituzionali verificati di adozione di sistemi di IA on-premise nella Pubblica Amministrazione, con l'obiettivo di identificare best practice, modelli replicabili e lezioni apprese.

% --- Nota metodologica ---
La metodologia adottata si basa su un approccio di analisi documentale e comparativa, con riferimento a fonti normative, istituzionali, accademiche e a report di organismi europei e nazionali. I casi studio sono stati selezionati sulla base della verificabilità delle fonti e della rilevanza per il contesto della PA italiana ed europea.
